\chapter{Triggers and filters}

Triggers decide when a workflow runs.
Filters decide when it should \emph{not} run.

\section{Core triggers}
\begin{itemize}[nosep]
  \item \code{push}
  \item \code{pull_request}
  \item \code{workflow_dispatch} (manual)
  \item \code{schedule} (cron)
  \item \code{workflow_call} (reusable workflows)
\end{itemize}

\section{Branch and path filters}
Run only when it matters.

\begin{lstlisting}[language=YAML]
on:
  push:
    branches: ["main"]
    paths:
      - "src/**"
      - ".github/workflows/**"
\end{lstlisting}

\section{PR triggers: common pitfalls}
\begin{itemize}[nosep]
  \item \code{pull_request} runs on PR events, not on merge commits.
  \item Use \code{pull_request_target} carefully; it runs with base repo privileges.
\end{itemize}

\begin{warningbox}{Warning: pull\_request\_target}
It is useful for labeling and triage, but it can expose secrets if you run untrusted code.
If you use it, do not checkout and run PR code with write permissions.
\end{warningbox}

\section{Manual inputs}
Manual runs are good for releases and maintenance.

\begin{lstlisting}[language=YAML]
on:
  workflow_dispatch:
    inputs:
      environment:
        description: "Target environment"
        required: true
        default: "staging"
\end{lstlisting}
